% GENERATED FILE. Edit canonical lessons, then run npm run generate:books.
% canonical-source-hash: fnv1a64:a67b411336fd5335
% canonical-lessons: LA-C329-mors, LA-C329-mortuus, LA-C329-exspiro, LA-C329-obeo, LA-C329-decedo

\chapter{Mors}
\label{ch:la-mors}
\hlchaptermodality{329}

\begin{chapteropening}
\textbf{By the end of this chapter:} \emph{I can talk about a death: death, dead, to breathe one's last, to die and to pass away.}

Complete the last lesson of chapter 329: I can talk about a death: death, dead, to breathe one's last, to die and to pass away.
\end{chapteropening}

\section[mors, mortis]{mors, mortis — death}

\label{lesson:LA-C329-mors}



\begin{quote}
Before the new one: say the Latin for a married woman, then the Latin for a divorce.
\end{quote}

\subsection*{You'll want to know: mors, mortis}
\textbf{mors, mortis} — \textquotedblleft{}death\textquotedblright{}.

\begin{quote}
Mors avī tōtam familiam trīstem fēcit.
\end{quote}

\begin{quote}
\emph{Grandfather's death made the whole family sad.}
\end{quote}

Feminine, 3rd declension, from the root of \emph{morior} (taught). English \emph{mortal} comes from its family.

\begin{cousinweb}
\textbf{mors} is inherited from Proto-Indo-European \emph{*mértis}, \textquotedblleft{}death\textquotedblright{} (EtymDB). English \emph{mortal} comes from its family, as does \emph{mortgage}, a dead pledge, through French.
\end{cousinweb}

\begin{grammarlens}[title={What you've built}]
One more word for telling, explaining and describing at B1.
\end{grammarlens}

\subsection*{Your turn}
\begin{itemize}
\raggedright
  \item \emph{Say it:} \emph{mors, mortis}
  \item \emph{Say it:} \emph{mors, mortis}, once more
  \item \emph{Say it:} a sentence of your own with \emph{mors, mortis}
  \item \emph{Recall:} say the Latin for a married woman, then the Latin for a divorce, then say \emph{mors, mortis} again
\end{itemize}

\begin{tcolorbox}[breakable,colback=teal!4,colframe=teal!35!black,title={Before you move on}]
What does mors, mortis mean? (\textquotedblleft{}Death\textquotedblright{}.) Say it once more.
\end{tcolorbox}
\section[mortuus, mortua, mortuum]{mortuus, mortua, mortuum — dead}

\label{lesson:LA-C329-mortuus}



\begin{quote}
Before the new one: say the Latin for a divorce, then the Latin for death.
\end{quote}

\subsection*{You'll want to know: mortuus, mortua, mortuum}
\textbf{mortuus, mortua, mortuum} — \textquotedblleft{}dead\textquotedblright{}.

\begin{quote}
Canis senex in angulō mortuus iacēbat.
\end{quote}

\begin{quote}
\emph{The old dog lay dead in the corner.}
\end{quote}

The past participle of \emph{morior}: \emph{mortuus est}, he died. As a noun, \emph{mortuus} is a dead man.

\begin{cousinweb}
\textbf{mortuus} is the participle of \emph{morior}, \textquotedblleft{}to die\textquotedblright{} (EtymDB), and goes back to Proto-Italic \emph{*mortwos}, \textquotedblleft{}dead\textquotedblright{}; Greek \emph{brotos}, \textquotedblleft{}mortal\textquotedblright{}, is a cousin.
\end{cousinweb}

\begin{grammarlens}[title={What you've built}]
One more word for telling, explaining and describing at B1.
\end{grammarlens}

\subsection*{Your turn}
\begin{itemize}
\raggedright
  \item \emph{Say it:} \emph{mortuus, mortua, mortuum}
  \item \emph{Say it:} \emph{mortuus, mortua, mortuum}, once more
  \item \emph{Say it:} a sentence of your own with \emph{mortuus, mortua, mortuum}
  \item \emph{Recall:} say the Latin for a divorce, then the Latin for death, then say \emph{mortuus, mortua, mortuum} again
\end{itemize}

\begin{tcolorbox}[breakable,colback=teal!4,colframe=teal!35!black,title={Before you move on}]
What does mortuus, mortua, mortuum mean? (\textquotedblleft{}Dead\textquotedblright{}.) Say it once more.
\end{tcolorbox}
\section[exspīrō, exspīrāre]{exspīrō, exspīrāre — to breathe one's last}

\label{lesson:LA-C329-exspiro}



\begin{quote}
Before the new one: say the Latin for death, then the Latin for dead.
\end{quote}

\subsection*{You'll want to know: exspīrō, exspīrāre}
\textbf{exspīrō, exspīrāre} — \textquotedblleft{}to breathe one's last\textquotedblright{}.

\begin{quote}
Senex in lectō exspīrāvit.
\end{quote}

\begin{quote}
\emph{The old man breathed his last in bed.}
\end{quote}

\emph{exspīrō, exspīrāre, exspīrāvī, exspīrātum}: \emph{ex} + \emph{spīrō} (taught), to breathe out. English \emph{expire}.

\begin{cousinweb}
\textbf{exspīrō} is built on \emph{spīrō}, \textquotedblleft{}to breathe\textquotedblright{}, with \emph{ex-}, \textquotedblleft{}out\textquotedblright{} (EtymDB): to breathe out the last breath. English \emph{expire} comes from it.
\end{cousinweb}

\begin{grammarlens}[title={What you've built}]
One more word for telling, explaining and describing at B1.
\end{grammarlens}

\subsection*{Your turn}
\begin{itemize}
\raggedright
  \item \emph{Say it:} \emph{exspīrō, exspīrāre}
  \item \emph{Say it:} \emph{exspīrō, exspīrāre}, once more
  \item \emph{Say it:} a sentence of your own with \emph{exspīrō, exspīrāre}
  \item \emph{Recall:} say the Latin for death, then the Latin for dead, then say \emph{exspīrō, exspīrāre} again
\end{itemize}

\begin{tcolorbox}[breakable,colback=teal!4,colframe=teal!35!black,title={Before you move on}]
What does exspīrō, exspīrāre mean? (\textquotedblleft{}To breathe one's last\textquotedblright{}.) Say it once more.
\end{tcolorbox}
\section[obeō, obīre]{obeō, obīre — to die; to go to meet}

\label{lesson:LA-C329-obeo}



\begin{quote}
Before the new one: say the Latin for dead, then the Latin for to breathe one's last.
\end{quote}

\subsection*{You'll want to know: obeō, obīre}
\textbf{obeō, obīre} — \textquotedblleft{}to die; to go to meet\textquotedblright{}.

\begin{quote}
Avia mea placidē in somnō obiit.
\end{quote}

\begin{quote}
\emph{My grandmother died peacefully in her sleep.}
\end{quote}

\emph{obeō, obīre, obiī, obitum}, like \emph{eō}. \emph{Mortem obīre}: to meet death; alone, \emph{obiit} means 'he died'. English \emph{obituary}.

\begin{cousinweb}
\textbf{obeō} is built on \emph{eō}, \textquotedblleft{}to go\textquotedblright{}, with \emph{ob}, \textquotedblleft{}towards\textquotedblright{} (EtymDB): to go to meet, and so to meet one's end. English \emph{obituary} comes from its noun \emph{obitus}, \textquotedblleft{}death\textquotedblright{}.
\end{cousinweb}

\begin{grammarlens}[title={What you've built}]
One more word for telling, explaining and describing at B1.
\end{grammarlens}

\subsection*{Your turn}
\begin{itemize}
\raggedright
  \item \emph{Say it:} \emph{obeō, obīre}
  \item \emph{Say it:} \emph{obeō, obīre}, once more
  \item \emph{Say it:} a sentence of your own with \emph{obeō, obīre}
  \item \emph{Recall:} say the Latin for dead, then the Latin for to breathe one's last, then say \emph{obeō, obīre} again
\end{itemize}

\begin{tcolorbox}[breakable,colback=teal!4,colframe=teal!35!black,title={Before you move on}]
What does obeō, obīre mean? (\textquotedblleft{}To die; to go to meet\textquotedblright{}.) Say it once more.
\end{tcolorbox}
\section[dēcēdō, dēcēdere]{dēcēdō, dēcēdere — to depart; to die}

\label{lesson:LA-C329-decedo}



\begin{quote}
Before the new one: say the Latin for to breathe one's last, then the Latin for to die.
\end{quote}

\subsection*{You'll want to know: dēcēdō, dēcēdere}
\textbf{dēcēdō, dēcēdere} — \textquotedblleft{}to depart; to die\textquotedblright{}.

\begin{quote}
Pater meus herī dēcessit.
\end{quote}

\begin{quote}
\emph{My father passed away yesterday.}
\end{quote}

\emph{dēcēdō, dēcēdere, dēcessī, dēcessum}: to go away, from a province or from life. English \emph{decease}.

\begin{cousinweb}
\textbf{dēcēdō} is built on \emph{cēdō}, \textquotedblleft{}to go, to withdraw\textquotedblright{}, with \emph{dē}, \textquotedblleft{}away\textquotedblright{} (Lewis \& Short). English \emph{decease} comes from its noun through French.
\end{cousinweb}

\begin{grammarlens}[title={What you've built}]
One more word for telling, explaining and describing at B1.
\end{grammarlens}

\subsection*{Your turn}
\begin{itemize}
\raggedright
  \item \emph{Say it:} \emph{dēcēdō, dēcēdere}
  \item \emph{Say it:} \emph{dēcēdō, dēcēdere}, once more
  \item \emph{Say it:} a sentence of your own with \emph{dēcēdō, dēcēdere}
  \item \emph{Recall:} say the Latin for to breathe one's last, then the Latin for to die, then say \emph{dēcēdō, dēcēdere} again
\end{itemize}

\begin{tcolorbox}[breakable,colback=teal!4,colframe=teal!35!black,title={Before you move on}]
What does dēcēdō, dēcēdere mean? (\textquotedblleft{}To depart; to die\textquotedblright{}.) Say it once more.
\end{tcolorbox}
